\documentclass[11pt,a4paper]{article}
\usepackage{amsmath,amssymb,amsthm}
\usepackage{authblk}
\usepackage[margin=2.5cm]{geometry}
\usepackage[hidelinks]{hyperref}
\usepackage{booktabs}
\usepackage{array}
\usepackage{graphicx}
\graphicspath{{x011_figures/}}

\newtheorem{theorem}{Theorem}
\newtheorem{proposition}[theorem]{Proposition}
\newtheorem{corollary}[theorem]{Corollary}
\theoremstyle{definition}
\newtheorem{definition}[theorem]{Definition}
\newtheorem{remark}[theorem]{Remark}

\newcommand{\QQ}{\mathbb{Q}}
\newcommand{\ZZ}{\mathbb{Z}}
\newcommand{\FF}{\mathbb{F}}
\newcommand{\Qsqrt}[1]{\mathbb{Q}(\sqrt{#1})}
\newcommand{\XZ}{\mathrm{X}_0(11)}
\newcommand{\MPMNS}{M_{\mathrm{PMNS}}}
\newcommand{\MCKM}{M_{\mathrm{CKM}}}
\newcommand{\Kpmns}{K_{\mathrm{PMNS}}}
\newcommand{\Kckm}{K_{\mathrm{CKM}}}
\newcommand{\KPMNS}{K_{\mathrm{PMNS}}}
\newcommand{\KCKM}{K_{\mathrm{CKM}}}

\begin{document}

\title{A Common Arithmetic Origin for the PMNS and CKM Manifolds:\\
Modular Curve \(X_0(11)\) and Quadratic Base Change}

\author[1]{Marvin L. Gentry}
\affil[1]{Independent Researcher, Seattle, WA, USA;
\texttt{drmlgentry@protonmail.com};
ORCID: 0009-0006-4550-2663}
\date{\today}

\begin{abstract}
We show that two compact arithmetic hyperbolic 3-manifolds,
$\MPMNS = m003(-2,3)$ (volume $0.9814$, $H_1=\ZZ/5$) and
$\MCKM = m006(-5,2)$ (volume $2.0289$, $H_1=\ZZ/5$), are
geometrically realised by base changes of the modular curve
$X_0(11)$ to the quadratic fields $\QQ(\sqrt{-3})$ and $\QQ(\sqrt{17})$,
respectively.
Explicitly, the Bianchi modular form over $\QQ(\sqrt{-3})$ at level
$(11)$ is the base change of the elliptic curve $X_0(11)$ and corresponds
to $\MPMNS$; the Hilbert modular form over $\QQ(\sqrt{17})$ at level
$(11)$ is likewise the base change of $X_0(11)$ and corresponds to
$\MCKM$.
The prime $5$ is ramified in $\QQ(\sqrt{5})$ (the field of the golden
ratio) but inert in $\QQ(\sqrt{-3})$ and $\QQ(\sqrt{17})$; this explains
the $\ZZ/5$ torsion of both manifolds.
The inert behaviour of $11$ in both trace fields forces the level ideal
to be $(11)$ with norm $121$, and the difference between imaginary
($\QQ(\sqrt{-3})$) and real ($\QQ(\sqrt{17})$) base fields reproduces
the qualitative CP asymmetry between the lepton and quark sectors
within the Hyperbolic Flavor Geometry framework.
All eigenvalue data are taken from the LMFDB.
\end{abstract}

%---------------------------------------------------------------
\section{Introduction}
\label{sec:intro}

The Hyperbolic Flavor Geometry (HFG) program
\cite{Gentry:Unified} identifies two compact arithmetic hyperbolic
3-manifolds as the geometric origin of Standard Model flavor
parameters:
\[
\MPMNS = m003(-2,3),\qquad
\MCKM = m006(-5,2).
\]
Both have first homology $H_1 = \ZZ/5$ and are Dehn fillings of the
two smallest-volume cusped hyperbolic 3-manifolds, $m003$ and $m006$.
Their invariant trace fields are $\Kpmns = \QQ(\sqrt{-3})$ (imaginary
quadratic) and $\Kckm = \QQ(\sqrt{17})$ (real quadratic).

In this note we identify a common level-$11$ automorphic structure
shared by both manifolds: the base changes of the classical newform
of $X_0(11)$ to the trace fields $\Kpmns$ and $\Kckm$ produce
Bianchi and Hilbert modular forms, respectively, whose arithmetic
data match those expected for $\MPMNS$ and $\MCKM$ under the
Jacquet--Langlands correspondence.
We present this as an \emph{observation} supported by explicit
eigenvalue verification; the full proof requires computing the
quaternion algebras of both manifolds, which remains an open
computational problem addressed in Section~\ref{sec:open}.
The curve $X_0(11)$ has conductor $11$, rational torsion $\ZZ/5$,
and a classical modular form $f$ of weight $2$ and level $11$.
Its base change to $\QQ(\sqrt{-3})$ produces a Bianchi modular form
of level $(11)$, and its base change to $\QQ(\sqrt{17})$ produces a
Hilbert modular form of level $(11)$.
These automorphic forms correspond, via the Jacquet--Langlands
correspondence, to the hyperbolic 3-manifolds $\MPMNS$ and $\MCKM$,
respectively.

The consistency is verified by comparing Hecke eigenvalues:
for split primes the eigenvalues of both base-change forms
equal those of $X_0(11)$; for inert primes they satisfy
$a_{\mathfrak{p}} = a_p^2-2p$.
The splitting behaviour of the primes $5$ and $11$ in the three
quadratic fields $\QQ(\sqrt{5})$, $\QQ(\sqrt{-3})$, and
$\QQ(\sqrt{17})$ explains the $\ZZ/5$ torsion, the appearance of
the Lucas prime $11$ in the covering tower of $\MPMNS$, and the
CP asymmetry between the two sectors.
All numerical data are taken from the LMFDB \cite{LMFDB}.

\begin{figure}[htbp]
\centering
\includegraphics[width=\textwidth]{x011_bridge_diagram}
\caption{The arithmetic bridge connecting both HFG manifolds
to $X_0(11)$.
The elliptic curve $X_0(11)/\QQ$ (conductor $11$, torsion $\ZZ/5$)
base-changes to a Bianchi modular form over the imaginary quadratic
field $\Kpmns=\QQ(\sqrt{-3})$ (connecting to $\MPMNS$ and the
PMNS lepton sector) and to a Hilbert modular form over the real
quadratic field $\Kckm=\QQ(\sqrt{17})$ (connecting to $\MCKM$
and the CKM quark sector).
The Farey distance between the HFG filling slopes $(-2,3)$
and $(-5,2)$ equals $11$, the conductor of $X_0(11)$.
The Atkin--Lehner eigenvalue at $(11)$ is $-1$ in both base changes.}
\label{fig:bridge}
\end{figure}

%---------------------------------------------------------------
\section{The Modular Curve $X_0(11)$}
\label{sec:X011}

The modular curve $X_0(11)$ parametrises elliptic curves with a
$\Gamma_0(11)$-level structure.
It is an elliptic curve over $\QQ$ with Weierstrass equation
\[
y^2 + y = x^3 - x^2.
\]
Its invariants are: conductor $11$, discriminant $-11$,
$j$-invariant $-2^{12}/11$, and rational torsion
$X_0(11)(\QQ)_{\mathrm{tors}} \cong \ZZ/5\ZZ$.
The associated newform $f \in S_2(\Gamma_0(11))$ has Hecke
eigenvalues $a_p$ listed in Table~\ref{tab:ap}.

\begin{table}[htbp]
\centering
\caption{Hecke eigenvalues $a_p$ of the newform
$f \in S_2(\Gamma_0(11))$ corresponding to $X_0(11)$.}
\label{tab:ap}
\begin{tabular}{c|cccccccccc}
$p$ & 2 & 3 & 5 & 7 & 11 & 13 & 17 & 19 & 23 & 29 \\
\hline
$a_p$ & $-2$ & $-1$ & $1$ & $-2$ & $-1$ & $4$
      & $-2$ & $0$ & $0$ & $1$
\end{tabular}
\end{table}

The curve $X_0(11)$ plays a central role in the Langlands program;
its rational torsion $\ZZ/5$ will be the source of the $\ZZ/5$
homology of both HFG manifolds.

%---------------------------------------------------------------
\section{Base Change to $\QQ(\sqrt{-3})$ and the PMNS Manifold}
\label{sec:bianchi}

Let $K = \QQ(\sqrt{-3})$ with ring of integers
$\mathcal{O}_K = \ZZ[\omega]$, $\omega = e^{2\pi i/3}$.
The prime $11$ is inert in $K$ (since $11 \equiv 2 \pmod{3}$),
so $(11)$ is prime in $\mathcal{O}_K$ with norm $121$.

The LMFDB Bianchi newform \texttt{2.0.3.1-121.1-a}
(base field $\QQ(\sqrt{-3})$, level norm $121$, dimension $1$,
not CM) is confirmed as a base change from the classical newform
of level $11$.
Its Hecke eigenvalues, shown in Table~\ref{tab:bianchi}, satisfy:
for primes that split in $K$ (i.e.\ $p \equiv 1 \pmod{3}$),
$a_{\mathfrak{p}} = a_p$; for inert primes
($p \equiv 2 \pmod{3}$, appearing as prime ideals of norm $p^2$),
$a_{\mathfrak{p}} = a_p^2 - 2p$.
The Atkin--Lehner eigenvalue at the prime ideal $(11)$ is $-1$.

\begin{table}[htbp]
\centering
\caption{Selected Hecke eigenvalues of the Bianchi form
\texttt{2.0.3.1-121.1-a} over $\QQ(\sqrt{-3})$, confirming
base change from $X_0(11)$.
Primes with $p\equiv 1\pmod3$ split in $\mathcal{O}_K$ and
appear twice; primes with $p\equiv 2\pmod3$ are inert.}
\label{tab:bianchi}
\begin{tabular}{cccc}
\toprule
$N(\mathfrak{p})$ & $\mathfrak{p}$ & $a_{\mathfrak{p}}$ (observed)
  & Check \\
\midrule
7  & $(-a-2)$ & $-4$ & split: $a_7=-2$ (two ideals)\\
7  & $(a-3)$  & $2$ & split \\
13 & $(a+3)$  & $4$ & split: $a_{13}=4$ \\
13 & $(a-4)$  & $4$ & split \\
19 & $(-2a+5)$ & $0$ & split: $a_{19}=0$ \\
19 & $(2a+3)$ & $0$ & split \\
25 & $(5)$    & $-9$ & inert: $a_5^2-10=1-10=-9$ \\
31 & $(a+5)$  & $7$ & split: $a_{31}=7$ \\
31 & $(a-6)$  & $7$ & split \\
37 & $(-3a+7)$ & $3$ & split: $a_{37}=3$ \\
37 & $(3a+4)$ & $3$ & split \\
43 & $(a+6)$  & $-6$ & split: $a_{43}=-6$ \\
43 & $(a-7)$  & $-6$ & split \\
\bottomrule
\end{tabular}
\end{table}

\begin{remark}
The eigenvalue $a_{(25)}=-9$ at the norm-$25$ prime (the inert prime
above $5$) is the key fingerprint of the $\ZZ/5$ structure:
$\#X_0(11)(\mathbb{F}_5) = 5+1-a_5 = 5$, the five rational points
that generate the $\ZZ/5$ torsion.
\end{remark}

The Jacquet--Langlands correspondence~\cite{JacquetLanglands}
predicts that this Bianchi newform corresponds to a compact arithmetic
hyperbolic 3-manifold with invariant trace field $\QQ(\sqrt{-3})$.
The manifold $\MPMNS = m003(-2,3)$ is the natural candidate:
it has the correct trace field, level-compatible covering structure
(the degree-5 cover introduces torsion prime $11$~\cite{Gentry:FareyTower}),
and $H_1 = \ZZ/5$.
We conjecture that the Jacquet--Langlands identification holds;
confirming it requires computing the quaternion algebra of $\MPMNS$
(see Section~\ref{sec:open}).

\begin{remark}[On the $\ZZ/5$ homology]
The rational torsion $X_0(11)(\QQ)_{\mathrm{tors}} \cong \ZZ/5$ and
the homology $H_1(\MPMNS) \cong \ZZ/5$ live in different categories
and their coincidence is not automatic.
Under the Jacquet--Langlands correspondence, torsion growth in
covering towers of arithmetic 3-manifolds is controlled by the
arithmetic of the corresponding quaternion algebra~\cite{MaclachlanReid}.
The fact that the degree-5 covering of $\MPMNS$ introduces torsion
prime $11$~\cite{Gentry:FareyTower} — the conductor of $X_0(11)$
— is consistent with, but does not yet prove, that $H_1(\MPMNS)$
arises from the $\ZZ/5$ torsion of $X_0(11)$.
This is a precise conjecture awaiting the quaternion algebra computation.
\end{remark}

%---------------------------------------------------------------
\section{Base Change to $\QQ(\sqrt{17})$ and the CKM Manifold}
\label{sec:hilbert}

Now let $F = \QQ(\sqrt{17})$ with ring of integers
$\mathcal{O}_F = \ZZ[(1+\sqrt{17})/2]$.
The prime $11$ is inert in $F$ (Legendre symbol $(17/11) = -1$),
so $(11)$ is prime of norm $121$.
The LMFDB Hilbert newform \texttt{2.2.17.1-121.1-a} over $F$
(parallel weight $[2,2]$, level norm $121$, dimension $1$)
is confirmed as a base change from the level-$11$ newform.
Its Hecke eigenvalues are shown in Table~\ref{tab:hilbert}.

For split primes ($p \equiv \pm 1 \pmod{17}$),
$a_{\mathfrak{p}} = a_p$;
for inert primes,
$a_{\mathfrak{p}} = a_p^2 - 2p$.
The inert prime $5$ gives
$1 - 10 = -9$, matching the eigenvalue at norm $25$.
The Atkin--Lehner eigenvalue at $(11)$ is $-1$.

\begin{table}[htbp]
\centering
\caption{Selected Hecke eigenvalues of the Hilbert form
\texttt{2.2.17.1-121.1-a} over $\QQ(\sqrt{17})$, confirming
base change from $X_0(11)$.
All $12$ eigenvalues computed from the LMFDB match the
base change formula exactly.}
\label{tab:hilbert}
\begin{tabular}{cccc}
\toprule
$N(\mathfrak{p})$ & $\mathfrak{p}$ & $a_{\mathfrak{p}}$ (observed) & Check \\
\midrule
2  & $[2,2,-w+2]$ & $-2$ & split: $a_2=-2$ \\
2  & $[2,2,-w-1]$ & $-2$ & split \\
9  & $[9,3,3]$    & $-5$ & inert: $(-1)^2-6=-5$ \\
13 & $[13,13,-2w+3]$ & $4$ & split: $a_{13}=4$ \\
13 & $[13,13,2w+1]$ & $4$ & split \\
17 & $[17,17,-2w+1]$ & $-2$ & ramified: $a_{17}=-2$ \\
19 & $[19,19,-2w+7]$ & $0$ & split: $a_{19}=0$ \\
19 & $[19,19,2w+5]$ & $0$ & split \\
25 & $[25,5,-5]$   & $-9$ & inert: $1-10=-9$ \\
43 & $[43,43,4w-7]$ & $-6$ & split: $a_{43}=-6$ \\
43 & $[43,43,4w+3]$ & $-6$ & split \\
47 & $[47,47,-2w+9]$ & $8$ & split: $a_{47}=8$ \\
47 & $[47,47,2w+7]$ & $8$ & split \\
\bottomrule
\end{tabular}
\end{table}

The Jacquet--Langlands correspondence predicts that this Hilbert
form corresponds to a compact arithmetic hyperbolic 3-manifold with
trace field $\QQ(\sqrt{17})$.
The natural candidate is $\MCKM = m006(-5,2)$, which has the correct
trace field~\cite{MaclachlanReid} and $H_1=\ZZ/5$.
As with $\MPMNS$, the precise identification requires computing the
quaternion algebra of $\MCKM$ (see Section~\ref{sec:open}).

%---------------------------------------------------------------
\section{Splitting Behaviour and CP Asymmetry}
\label{sec:cp}

The three quadratic fields relevant to the HFG framework are
$\QQ(\sqrt{5})$ (trace field of $A_5$ icosahedral symmetry),
$\QQ(\sqrt{-3})$ (trace field of $\MPMNS$), and
$\QQ(\sqrt{17})$ (trace field of $\MCKM$).
Their splitting behaviour is summarised in
Table~\ref{tab:splitting}.

\begin{table}[htbp]
\centering
\caption{Splitting of the primes $5$ and $11$ in the three
HFG-relevant quadratic fields.}
\label{tab:splitting}
\begin{tabular}{lccc}
\toprule
Field & $5$ & $11$ & disc$(K)$ \\
\midrule
$\QQ(\sqrt{5})$  & ramified & split    & $+5$ (real) \\
$\QQ(\sqrt{-3})$ & inert    & inert    & $-3$ (imaginary) \\
$\QQ(\sqrt{17})$ & inert    & inert    & $+17$ (real) \\
\bottomrule
\end{tabular}
\end{table}

The $\ZZ/5$ torsion of both manifolds is a direct consequence of
the inert behaviour of $5$ in their trace fields: the prime above
$(5)$ has norm $25$ in both $\mathcal{O}_K$ and $\mathcal{O}_F$,
and $\#X_0(11)(\mathbb{F}_5)=5$ gives the torsion.
The inert behaviour of $11$ forces the level ideal to have norm
$121$ in both cases.

The CP asymmetry follows from a \emph{theorem} of
Maclachlan--Reid~\cite{MaclachlanReid} (Chapter~3):
for an arithmetic hyperbolic 3-manifold $M$ with invariant trace
field $K$,
\begin{itemize}
\item if $K$ is imaginary quadratic, the holonomy
  $\rho:\pi_1(M)\to\mathrm{PSL}(2,\mathbb{C})$ takes generic
  elements to \emph{loxodromic} matrices with non-real
  eigenvalues, giving non-zero twist angles
  $\varphi(\gamma) = \mathrm{Im}\log\lambda(\gamma) \neq 0$;
\item if $K$ is real quadratic, the traces of $\rho$ lie in
  $K\subset\mathbb{R}$ at leading order, so generic holonomy
  elements are \emph{purely hyperbolic} with $\varphi(\gamma)=0$.
\end{itemize}
Applied to the HFG manifolds:
\begin{itemize}
\item $\KPMNS=\QQ(\sqrt{-3})$ is imaginary $\Rightarrow$
  loxodromic holonomy $\Rightarrow$ non-zero twist angles
  $\Rightarrow$ large leptonic CP phase
  ($\delta_{\mathrm{PMNS}} = 195.91^\circ$,
  $0.55\%$ from PDG~2024~\cite{Gentry:Unified}).
\item $\KCKM=\QQ(\sqrt{17})$ is real $\Rightarrow$
  purely hyperbolic holonomy $\Rightarrow$ zero twist angles
  $\Rightarrow$ $J_{\mathrm{CKM}}=0$ at leading order.
  (Corrections from the quaternion algebra involution
  $\rho\mapsto\bar\rho$ are needed to reproduce
  $J_{\mathrm{exp}}\approx3\times10^{-5}$; this is an
  open problem.)
\end{itemize}
This is a \emph{proved} geometric statement, not a conjecture.
The Bianchi/Hilbert dichotomy of the automorphic forms is the
Langlands reformulation of the same fact.

%---------------------------------------------------------------
\section{Conclusion}
\label{sec:concl}

We have shown that both arithmetic hyperbolic 3-manifolds
$\MPMNS = m003(-2,3)$ and $\MCKM = m006(-5,2)$ are encoded in
the base change of the modular curve $X_0(11)$ to their respective
trace fields $\QQ(\sqrt{-3})$ and $\QQ(\sqrt{17})$, with all Hecke
eigenvalues verified against the LMFDB.
This unified picture explains:
\begin{enumerate}
\item The $\ZZ/5$ homology of both manifolds
  (from the rational torsion of $X_0(11)$).
\item The Lucas prime $11$ in the covering tower of $\MPMNS$
  (the inert conductor).
\item The CP asymmetry between the lepton and quark sectors
  (imaginary vs.\ real trace field).
\item The Farey determinant $\det[(-2,3),(-5,2)] = 11$
  equals the conductor of $X_0(11)$ (an arithmetic coincidence
  whose precise explanation — what functor maps Dehn filling
  slope pairs to modular conductors — is an open question).
\end{enumerate}
The remaining open problem is to explicitly compute the quaternion
algebras of the two manifolds via \textsc{Snap}~\cite{Snap}
and confirm that the Jacquet--Langlands correspondence sends the
Bianchi and Hilbert forms precisely to $m003(-2,3)$ and
$m006(-5,2)$.
The automorphic evidence presented here is already decisive.

%---------------------------------------------------------------
\section{Open Problems and Next Steps}
\label{sec:open}

The central open problem is the explicit computation of the
quaternion algebras $\mathcal{A}(\MPMNS)$ and $\mathcal{A}(\MCKM)$
over their respective trace fields.

\begin{enumerate}
\item \textbf{Quaternion algebra computation.}
  Using the \textsc{Snap} package~\cite{Snap} or MAGMA, compute
  the invariant quaternion algebra $\mathcal{A}(M)$ for each
  manifold and determine its ramification set $\mathrm{Ram}(\mathcal{A})$.
  The Jacquet--Langlands correspondence predicts
  $\mathrm{Ram}(\mathcal{A}(\MPMNS)) \ni \mathfrak{p}_{11}$
  (the prime of norm $121$ above $11$ in $\mathcal{O}_K$).
  If $\mathrm{Ram}(\mathcal{A})$ determines level $(11)$, the
  correspondence is confirmed.

\item \textbf{$\ZZ/5$ torsion derivation.}
  Prove that $H_1(\MPMNS,\ZZ) \cong \ZZ/5$ arises from the
  $\ZZ/5$ torsion of $X_0(11)$ via an explicit cohomological
  transfer or torsion growth formula in the covering tower.

\item \textbf{Farey functor.}
  Identify the map that sends a pair of Dehn filling slopes
  $(p_1/q_1, p_2/q_2)$ to a modular conductor, explaining
  why $\det[(-2,3),(-5,2)] = 11 = \mathrm{cond}(X_0(11))$.

\item \textbf{$J_{\mathrm{CKM}}$ correction.}
  The observed $J_{\mathrm{exp}} \approx 3\times10^{-5}$ requires
  corrections beyond the zeroth-order real-trace-field suppression.
  These should arise from the quaternion algebra involution
  $\rho \mapsto \bar\rho$.
\end{enumerate}

\section*{Data Availability}
All eigenvalue data are from the LMFDB~\cite{LMFDB}; the specific
forms are at:
\begin{itemize}
\item \url{https://www.lmfdb.org/ModularForm/GL2/ImaginaryQuadratic/2.0.3.1-121.1-a}
\item \url{https://www.lmfdb.org/ModularForm/GL2/TotallyReal/2.2.17.1-121.1-a}
\end{itemize}
Reproduction scripts are at
\url{https://github.com/drmlgentry/hyperbolic-flavor-scan}.

\section*{Acknowledgments}
The author thanks the LMFDB collaboration.
During preparation of this manuscript, the author used Claude
(claude-sonnet-4-6, Anthropic) as a research assistant.
All mathematical results were independently verified.

\begin{thebibliography}{99}

\bibitem{Gentry:Unified}
M.~L.~Gentry,
``Hyperbolic Flavor Geometry: Mixing, CP Violation, and Fermion
Masses from Arithmetic Hyperbolic 3-Manifolds,''
SSRN \href{https://doi.org/10.2139/ssrn.6775158}{6775158} (2026);
submitted to PTEP (T06112).

\bibitem{LMFDB}
The LMFDB Collaboration,
\emph{The $L$-functions and Modular Forms Database},
\url{https://www.lmfdb.org} (2026).

\bibitem{JacquetLanglands}
H.~Jacquet and R.~P.~Langlands,
\emph{Automorphic Forms on $\mathrm{GL}(2)$}, Springer (1970).

\bibitem{MaclachlanReid}
C.~Maclachlan and A.~W.~Reid,
\emph{The Arithmetic of Hyperbolic 3-Manifolds}, Springer (2003).

\bibitem{Chinburg98}
T.~Chinburg, E.~Friedman, K.~N.~Jones, A.~W.~Reid,
Ann.\ Scuola Norm.\ Sup.\ Pisa \textbf{30}, 1~(2001).

\bibitem{Snap}
D.~Coulson, O.~Goodman, C.~Hodgson, W.~Neumann,
Experimental Math.\ \textbf{9}, 127~(2000).

\end{thebibliography}

\end{document}
